% Sec. V (Verification of the Model): V-A The Core Model and V-B Verification Results.
% Rules described here in prose are printed in sections/appx_F3.tex (Appendix B).
\section{Verification of the Model}\label{sec:verification}

We describe the model (Sec.~\ref{sec:core}), what we verified (Sec.~\ref{sec:verified}) and the
resulting guarantees (Secs.~\ref{sec:properties} to~\ref{sec:wormhole}).
Listings are verbatim; Appendix~\ref{app:listings} prints further rules.

\subsection{The Core Model}\label{sec:core}

\subsubsection{Payment Channels}\label{sec:core-channels}
A channel is identified by a fresh pointer \code{ptr} to its founding transaction. Our model
tracks two mirror facts about this channel, \code{StateA} and \code{StateB}, one held by each
party, each carrying the two balances (\code{fA}, \code{fB}), a state index~\code{n}, and that
party's own revocation secret together with the corresponding commitment. The commitment is
stored without its signature; signatures are checked when a state is installed
(\code{Both_States_Committed} for the first, \code{Revoke_Old_Secret} for every later one) as
\code{Eq} action facts, enforced by restriction \code{Equality}.

The heart of the model is the \emph{update dance} (Appendix~\ref{lst:update}): commit a new
state first, then revoke the old one, so that the old state is never revoked until a fully
signed new one exists. \code{Update_State} samples the next state index and both fresh
revocation secrets, reads the new balances from the environment, and turns the two mirrors into
two pending mirrors that keep the old state next to the new, signed one. Only afterwards does
\code{Revoke_Old_Secret} install the new state, recording the old revocation secrets as
persistent knowledge (\code{!RevokedSecret}) and the old commitments as linear one-shot tokens
(\code{RevokedCommitmentUnspent}, Sec.~\ref{par:single-spend}).

\code{Lock_Funds_And_Open} mints the two mirrors, the persistent facts
\code{!ValidOnChain(ptr)}, \code{!ChannelRecord(ptr,$A,$B)} and \code{!ChannelConnect} (one per
direction), and the per-channel slot \code{Mutex(ptr)}. The remaining phases reuse these facts
(Appendix~\ref{lst:close}). A unilateral close spends the current mirror:
\code{Publish_Current_A} turns \code{StateA} into an \code{OnChainDelayed} and an
\code{OnChainInstant} output, and \code{Claim_Funds_A} later releases the delayed output to its
owner; its action \code{CSVPassed} captures that time has passed. A cooperative close
(\code{Honest_Close}) signs one shared settlement transaction instead, and punishment
(\code{Publish_Revoked_A}, then \code{Punish_A}) transfers the whole channel to the honest party.

\subsubsection{Routing}\label{sec:core-routing}
Each direction of the routing layer (Sec.~\ref{sec:routing}) has its own fact. HTLC
offers travel forward as linear facts \code{Route($Prev,$Me,ptrIn,y)}, each also recorded in a
persistent fact \code{!RouteInfo}; once an HTLC is redeemed, the preimage travels back as
\code{Fulfilled(ptr,x)}. One generic forwarding rule serves any number of hops:
\begin{spthy}
rule Forward_HTLC:
    [ Route($Prev, $Me, ptrIn, y),
      !ChannelConnect(ptrOut, $Me, $Next),
      Mutex(ptrOut),
      In(%vIn), In(%fee), In(%vOut) ]
    --[ HTLCForwarded($Me, $Prev, $Next, ptrIn, ptrOut, y),
        HTLCAdded(ptrOut, y),
        FeeCharged($Me, y, %fee),
        Eq(%vIn,  %vOut %+ %fee),
        Neq(ptrIn, ptrOut), Neq($Prev, $Me), Neq($Me, $Next) ]->
    [ IntermediaryPending($Me, $Prev, $Next, ptrIn, ptrOut, y, %vIn, %vOut, %fee),
      Route($Me, $Next, ptrOut, y),
      !RouteInfo($Me, $Next, ptrOut, y) ]
\end{spthy}
It consumes the incoming \code{Route} fact, produces the outgoing one, takes the slot of the
outgoing channel, and stores in \code{IntermediaryPending}
everything the hop needs to settle or be refunded later. \code{$Next} is any party with which
\code{$Me} shares a channel, so routes are built hop by hop and have no fixed length.
\code{HTLCForwarded} identifies the hop in our lemmas, \code{HTLCAdded} feeds restriction
\code{OneHTLCPerHashPerChannel}, and \code{FeeCharged} records the fee the hop expects; the
\code{Neq} facts make the forwarder distinct from its neighbours and the two channels distinct.

The settle round runs backwards (Appendix~\ref{lst:settle}). \code{Receiver_Create_Invoice}
signs and sends a tagged invoice naming the receiver, the sender and the payment hash
$y = h(x)$ of a fresh preimage~$x$, kept private in \code{ReceiverInv};
\code{Sender_Offer_HTLC} checks the signature before offering the first HTLC.
\code{Receiver_Fulfill_HTLC} consumes the last \code{Route} fact, checks that $x$ hashes to~$y$,
emits \code{Redeemed(ptrIn,y)} and passes \code{Fulfilled(ptrIn,x)} upstream;
\code{Intermediary_Settle} turns \code{IntermediaryPending} and \code{Fulfilled(ptrOut,x)} into
\code{Fulfilled(ptrIn,x)}, recording \code{FeeEarned}. Both free \code{Mutex(ptrIn)};
\code{Sender_Settle} closes the payment (\code{SenderSettled}). Once \code{Open_Timeout} times a
channel out (an untimed event, Sec.~\ref{sec:threat}), \code{Sender_Timeout_Refund} and
\code{Intermediary_Timeout_Refund} refund its pending HTLC without punishment
(\code{Refunded}) and free the slot.

\subsubsection{The Wormhole Rule}\label{sec:core-skip}
Besides the Dolev-Yao network and key compromise (\code{Compromise_Ltk}), the adversary has one
routing rule:
\begin{spthy}
rule AttackerSkip:
    [ Fulfilled(ptrExit, x),
      !RouteInfo(ExitAtt, Next, ptrExit, y)[-,no_precomp],
      !RouteInfo(EntryAtt, Skipped, ptrSkip, y)[-,no_precomp],
      !Compromised(ExitAtt),
      !Compromised(EntryAtt) ]
    --[ FulfilledProduced(ptrSkip, x),
        Skipped(ptrSkip, x),
        Eq(h(x), y) ]->
    [ Fulfilled(ptrSkip, x) ]
\end{spthy}
If two compromised parties, \code{ExitAtt} and \code{EntryAtt}, have added HTLCs for the same
payment hash~$y$ on channels \code{ptrExit} and \code{ptrSkip}, \code{AttackerSkip} delivers the preimage released on \code{ptrExit} to \code{ptrSkip}. \code{EntryAtt} can then
settle upstream, while the hops between \code{ptrSkip} and \code{ptrExit}, starting with
\code{Skipped}, never see their outgoing HTLC fulfilled and can at best be refunded. The rule
emits \code{Skipped(ptrSkip,x)} but neither \code{Redeemed} nor \code{FeeEarned}; it yields the
wormhole trace of Sec.~\ref{sec:wormhole} (\code{Wormhole_Fee_Theft_Reachable}).

\subsubsection{Amounts and Fees}\label{sec:core-amounts}
Amounts and fees are modelled per hop only. \code{Forward_HTLC} reads the incoming amount, the
fee and the outgoing amount from the network as natural numbers and requires the incoming
amount to equal the outgoing amount plus the fee, so this balance holds at every hop by
construction. \code{Sender_Offer_HTLC}, the invoice and the \code{Route} facts carry no amount,
so amounts are not chained along the path, and the model says nothing about end-to-end
value, such as the amount the receiver obtains or the total fee the sender pays. We keep per-hop
fees because the wormhole is a fee theft: \code{FeeCharged} is emitted when a forwarder locks
its outgoing HTLC, \code{FeeEarned} only when it later claims its incoming one, and in the
wormhole trace the bypassed honest hop emits only the first, although the
sender settles and the receiver is paid. A hop whose key is never compromised earns a
fee only once its outgoing HTLC is redeemed (Lemma~\ref{lem:fee}). The fee is one per-hop
quantity, not Lightning's \texttt{base\_fee + fee\_rate} formula.

Three non-goals of Sec.~\ref{sec:structure} show in the rules. \code{Forward_HTLC} has no
\code{Wallet} premise and HTLCs do not change the balances in \code{StateA} and \code{StateB},
so a forward never fails for lack of funds. Route information is kept in local state, not in
layered encryption: the sender chooses only the first hop, each forwarder its next hop among its
own channels, and the wormhole appears purely as fee theft, as also characterised
in~\cite{malavolta2019}. The mirrors abstract the output structure of
Fig.~\ref{fig:transaction}, and the two-stage HTLC claim becomes a single transition; for the
channel balance, the second stage's CSV delay and revocation path are enforced by
\code{Publish_Current_A/B}, \code{Claim_Funds_A/B} and the punishment rules.

\subsection{Verification Results}\label{sec:verified}

\begin{table}[t]
  \caption{Verification runs of \texttt{lightning.spthy} (\tamarin{} 1.10), in the order of
    \texttt{make all}; \texttt{aux} runs \texttt{lightning\_aux.spthy} (Table~\ref{tab:aux}). A, E: all-traces and
    exists-trace lemmas.}
  \label{tab:results}
  \centering
  \footnotesize
  \setlength{\tabcolsep}{3pt}
  \begin{tabular}{@{}L{1.4cm}L{6.1cm}rr@{}}
    \toprule
    Target & File, flags and content & A & E\\
    \midrule
    \texttt{verify} & All restrictions: 32 safety, 13 hop-bypass, 10 exists-trace (1 replayed)
      & 45 & 10\\
    \texttt{general} & \code{-D=GENERAL}, \code{NO_*}: T1 to T3 removed; Theorem~\ref{thm:anyhop}
        (7 steps), 7 helpers, 16 \texttt{verify} lemmas re-proved
      & 24 & 0\\
    \texttt{aux} & \code{-D=PTLC}: countermeasure on a fixed route & 1 & 4\\
    & \code{-D=HTLC}: baseline, wormhole reachable & 0 & 3\\
    & \code{-D=CLOCK}: block clock & 6 & 3\\
    & \code{-D=CLTV}: CLTV-delta arithmetic & 3 & 0\\
    & \code{-D=TIMEOUT}: early-timeout race reachable & 0 & 1\\
    \texttt{necessity} & T3, T2 or T1 removed per run (Table~\ref{tab:necessity}): 5 liveness
        lemmas, 600\,s each: 10/15 verified, 5 without result; T3 counterexample (28 steps)
      & 5$\times$3 & 1\\
    \midrule
    \textbf{Total} & \texttt{verify}, \texttt{general}, \texttt{aux}: 100/100 verified
      & 79 & 21\\
    \bottomrule
  \end{tabular}
\end{table}

Table~\ref{tab:results} lists every verification run of our artifact, all obtained with \tamarin{} 1.10
(which reports the \maude{} version used, 3.2, as unsupported) and reproduced on a laptop with a
6\,GB heap limit (\texttt{MEM=6G}). The default mode (\texttt{make verify}) proves the safety
lemmas about channels and payments (Secs.~\ref{sec:properties} to~\ref{sec:wormhole},
Appendix~\ref{app:catalogue}) and that, if the sender settles, the forwarders at positions 1 to 6
of a route are bypassed only if some key was compromised (Sec.~\ref{sec:bypass-fixed}). Its exists-trace lemmas show that the model is not vacuous (channel updates, cooperative closes,
refunds and honest multi-hop payments are reachable), that the assumptions of our safety lemmas
matter (a payment without a matching invoice can settle once the receiver's key is compromised;
a compromised intermediary can lose), and that the wormhole is reachable. \texttt{make general}
drops T1 to T3 and proves the bypass theorem for routes of any length,
for payments with a single offer per payment hash (Sec.~\ref{sec:bypass-any}); since removing
restrictions only adds traces, this all-traces result also holds in the default mode.

% The label sits before \paragraph so that it prints as "V-B" (IEEEtran would print "V-B0a").
\phantomsection\label{par:single-spend}%
\paragraph{The Linear Single-Spend Discipline}
We enforce the no-double-spend rule using linear facts. Our first
\code{Publish_Revoked_A} lacked the linear premise \code{RevokedCommitmentUnspent}, so punishment was gated only by the persistent
\code{!RevokedSecret} and could fire again and again for the same revoked state: if a party
cheats, the honest party can take the channel funds many times. Lemmas like
\code{No_Punishment_Without_Cheating} still pass, because a cheat is still required; what they
miss is that punishment should happen only once. Now \code{Revoke_Old_Secret} also creates this
token for both parties, and \code{Publish_Revoked_A} (symmetrically
\code{Publish_Revoked_B}) must consume it, so a revoked commitment is punished only once. We
found the same critical flaw in our initial routing layer: with the persistent
\code{!ChannelConnect} alone, one channel could add unboundedly many HTLCs on the same output.
The linear slot \code{Mutex(ptr)} fixes it: produced once at opening, consumed by adding an
HTLC and given back when that HTLC is redeemed or refunded, it lets a channel carry at most one
pending HTLC at a time and be reused, each time with a payment hash not used on it before
(\code{OneHTLCPerHashPerChannel}). Adding a linear premise
can remove exists-trace witnesses, so we re-ran \code{lightning.spthy} after both
fixes: all 55 lemmas of \texttt{make verify} are verified, including the ten exists-trace lemmas.
